%================================================================
% LAPORAN AKHIR: SMART INSECT IDENTIFIER & AI INSIGHTS
% Mata Kuliah : Praktikum Pembelajaran Mesin
% Penulis     : Faris Zain As-Shadiq (2308107010039)
% Universitas : Universitas Syiah Kuala
%================================================================

\documentclass[12pt, a4paper]{article}

%----------------------------------------------------------------
% PACKAGES
%----------------------------------------------------------------
\usepackage[a4paper,
            top=2.5cm, bottom=2.5cm,
            left=3cm,  right=2.5cm]{geometry}
\usepackage[T1]{fontenc}
\usepackage[utf8]{inputenc}
\usepackage{times}            % Times New Roman
\usepackage{graphicx}
\usepackage{ulem}             % \uline
\usepackage{setspace}
\usepackage{parskip}          % spasi antar paragraf
\usepackage{indentfirst}      % indent paragraf pertama
\usepackage{booktabs}         % tabel profesional
\usepackage{tabularx}         % tabel lebar fleksibel
\usepackage{longtable}        % tabel multi-halaman
\usepackage{array}
\usepackage{multirow}
\usepackage[table,xcdraw]{xcolor}
\usepackage{fancyhdr}
\usepackage{titlesec}
\usepackage{caption}
\usepackage{float}
\usepackage{hyperref}
\usepackage{enumitem}
\usepackage{microtype}
\usepackage{tocloft}          % daftar isi kustom

%----------------------------------------------------------------
% HYPERREF SETUP
%----------------------------------------------------------------
\hypersetup{
    colorlinks = true,
    linkcolor  = black,
    urlcolor   = blue,
    citecolor  = black,
    pdftitle   = {Laporan Final Project Praktikum Pembelajaran Mesin},
    pdfauthor  = {Faris Zain As-Shadiq}
}

%----------------------------------------------------------------
% CONFIGURATION: TOCOFT (DAFTAR ISI, GAMBAR, & TABEL)
%----------------------------------------------------------------
\renewcommand{\contentsname}{DAFTAR ISI}
\renewcommand{\listfigurename}{DAFTAR GAMBAR}
\renewcommand{\listtablename}{DAFTAR TABEL}

\renewcommand{\cfttoctitlefont}{\hfill\large\bfseries}
\renewcommand{\cftaftertoctitle}{\hfill}
\setlength{\cftbeforetoctitleskip}{-10pt}

\renewcommand{\cftloftitlefont}{\hfill\large\bfseries}
\renewcommand{\cftafterloftitle}{\hfill}
\setlength{\cftbeforeloftitleskip}{-10pt}

\renewcommand{\cftlottitlefont}{\hfill\large\bfseries}
\renewcommand{\cftafterlottitle}{\hfill}
\setlength{\cftbeforelottitleskip}{-10pt}

%----------------------------------------------------------------
% HEADING STYLE (FORMAT BAB NORMAL INDONESIA)
%----------------------------------------------------------------
\newcommand{\sectionbreak}{\clearpage}

\titleformat{\section}[display]
  {\large\bfseries\centering}
  {BAB \thesection}
  {0pt}
  {\large\bfseries\MakeUppercase}
\titlespacing*{\section}{0pt}{-10pt}{18pt}

\titleformat{\subsection}
  {\normalsize\bfseries}
  {\thesubsection}{0.8em}{}

\titleformat{\subsubsection}
  {\normalsize\bfseries\itshape}
  {\thesubsubsection}{0.8em}{}

%----------------------------------------------------------------
% HEADER / FOOTER
%----------------------------------------------------------------
\pagestyle{fancy}
\fancyhf{} 
\fancyfoot[C]{\thepage} 
\renewcommand{\headrulewidth}{0pt} 
\renewcommand{\footrulewidth}{0pt}

%----------------------------------------------------------------
% CAPTION STYLE
%----------------------------------------------------------------
\captionsetup[figure]{labelfont=bf, font=small, justification=centering}
\captionsetup[table]{labelfont=bf, font=small, justification=centering}

%----------------------------------------------------------------
% WARNA TABEL
%----------------------------------------------------------------
\definecolor{tblhead}{RGB}{16,185,129}     % hijau emerald header
\definecolor{tblalt}{RGB}{240,253,250}    % hijau muda alt row
\definecolor{tblkey}{RGB}{204,251,241}    % hijau muda key col

%----------------------------------------------------------------
% SPACING & LISTINGS
%----------------------------------------------------------------
\onehalfspacing
\setlength{\parindent}{1.25cm}
\setlength{\parskip}{6pt}

\usepackage{listings}         % untuk source code
\lstset{
  basicstyle=\ttfamily\small,
  backgroundcolor=\color{gray!5},
  frame=single,
  breaklines=true,
  captionpos=b,
  numbers=left,
  numberstyle=\tiny\color{gray},
  keywordstyle=\color{blue},
  commentstyle=\color{green!50!black},
  stringstyle=\color{red!70!black},
  extendedchars=true,
  literate={│}{|}1 {├}{|}1 {─}{-}1 {└}{|}1,
}

%================================================================
\begin{document}
\setcounter{page}{0}

%================================================================
% HALAMAN SAMPUL
%================================================================
\begin{titlepage}
  \thispagestyle{empty}
  
  \noindent\textit{Tugas Final Project Praktikum Pembelajaran Mesin}

  \begin{center}

    \vspace{2.2cm}

    {\large\bfseries
      LAPORAN FINAL PROJECT \\[4pt]
      SMART INSECT IDENTIFIER \& AI INSIGHTS \\[4pt]
      SISTEM KLASIFIKASI SERANGGA BERBASIS DEEP LEARNING (LENSARTHROPODA)
    }

    \vspace{2.2cm}

    disusun untuk memenuhi tugas akhir mata kuliah \\
    Praktikum Pembelajaran Mesin

    \vspace{1.4cm}

    oleh:

    \vspace{0.3cm}

    \uline{Faris Zain As-Shadiq} \\
    2308107010039

    \vspace{1.8cm}

    \includegraphics[width=4cm]{Screenshot 2026-05-10 214133.png}

    \vspace{1.8cm}

    {\bfseries
      JURUSAN INFORMATIKA \\
      FAKULTAS MATEMATIKA DAN ILMU PENGETAHUAN ALAM \\
      UNIVERSITAS SYIAH KUALA \\
      2026
    }

  \end{center}
\end{titlepage}

%================================================================
% HALAMAN DEPAN
%================================================================
\newpage
\tableofcontents

\newpage
\listoffigures

\newpage
\listoftables

%================================================================
% BAB 1: PENDAHULUAN
%================================================================
\newpage
\section{Pendahuluan}

\subsection{Latar Belakang}
Sektor pertanian memegang peranan krusial dalam menyokong perekonomian dan ketahanan pangan global. Namun, produktivitas lahan pertanian seringkali terancam oleh gangguan hama tanaman yang sebagian besar merupakan spesies serangga. Identifikasi spesies serangga secara cepat dan akurat menjadi kunci utama dalam merumuskan strategi penanggulangan hama terpadu. Metode manual yang bergantung pada keahlian entomologi membutuhkan waktu lama, mahal, dan tidak efisien untuk skala lapangan luas. 

Perkembangan teknologi \textit{Deep Learning} khususnya \textit{Convolutional Neural Networks} (CNN) telah membuka peluang besar untuk mengotomatisasi proses pengenalan visual ini. Dengan mengombinasikan kekuatan klasifikasi citra dari arsitektur model \textit{EfficientNet-B3} dan wawasan biologis dinamis dari \textit{Large Language Model} (LLM) Gemini, sistem ini tidak hanya mampu mengenali nama serangga, melainkan juga menyajikan informasi ekosistem secara utuh bagi pengguna.

\subsection{Rumusan Masalah}
Adapun rumusan masalah yang diselesaikan dalam proyek pengembangan sistem ini adalah:
\begin{enumerate}
    \item Bagaimana cara memproses citra dataset serangga pertanian dengan keanekaragaman tinggi agar model dapat belajar secara optimal tanpa terhambat file gambar rusak?
    \item Bagaimana performa transfer learning menggunakan arsitektur \textit{EfficientNet-B3} untuk mengenali 118 kelas serangga pertanian?
    \item Bagaimana cara mengintegrasikan hasil prediksi model \textit{PyTorch} dengan Gemini API secara aman, andal, dan minim latensi?
\end{enumerate}

\subsection{Tujuan Proyek}
Tujuan utama proyek ini meliputi:
\begin{enumerate}
    \item Membangun alur prapemrosesan data tangguh yang menangani berkas citra korup secara otomatis.
    \item Melatih dan mengevaluasi model klasifikasi serangga berskala 118 kelas dengan target akurasi uji tinggi.
    \item Merancang antarmuka web modern dengan interaksi animasi premium yang menyatukan prediksi model lokal dan wawasan teks AI.
\end{enumerate}

%================================================================
% BAB 2: EVALUASI MODEL DEEP LEARNING
%================================================================
\newpage
\section{Evaluasi Model Deep Learning}

\subsection{Preprocessing & Augmentasi Data}
Dataset citra serangga pertanian yang digunakan dikelompokkan ke dalam tiga subset utama: \textit{train}, \textit{val}, dan \textit{test}. Salah satu kendala utama pada dataset besar adalah adanya citra korup yang dapat menghentikan proses latih. Penanganan dilakukan melalui pemindaian proaktif dengan blok \textit{try-except} memanfaatkan pustaka \textit{Pillow} (\textit{UnidentifiedImageError}). Citra yang lolos kemudian diproses dengan augmentasi berikut:
\begin{itemize}
    \item \textbf{Resize & Random Crop}: Citra diubah ukurannya ke dimensi $320 \times 320$ piksel kemudian dipotong acak menjadi $300 \times 300$ piksel untuk menyesuaikan resolusi input optimal \textit{EfficientNet-B3}.
    \item \textbf{Geometric Transformations}: Penerapan \textit{RandomHorizontalFlip} ($p=0.5$), \textit{RandomVerticalFlip} ($p=0.2$), dan \textit{RandomRotation} ($20^\circ$).
    \item \textbf{Color Transformations}: Penggunaan \textit{ColorJitter} (kecerahan, kontras, saturasi) untuk menyimulasikan perbedaan pencahayaan kamera di lapangan.
    \item \textbf{Normalisasi}: Penggunaan rata-rata dan deviasi standar standar ImageNet untuk mempercepat konvergensi laju pemelajaran.
\end{itemize}

\subsection{Arsitektur Model & Pelatihan}
Sistem ini menggunakan teknik \textit{Transfer Learning} berbasis \textbf{EfficientNet-B3} yang telah dilatih pada ImageNet. Penggunaan arsitektur ini didasarkan pada efisiensi penskalaan koefisien gabungan (\textit{compound scaling}) yang menyeimbangkan kedalaman, lebar, dan resolusi jaringan. Model dioptimasi dengan scheduler laju laju pemelajaran \textit{CosineAnnealingLR} laju awal $10^{-4}$ dan bobot peluruhan laju $10^{-4}$. Kerugian dihitung melalui \textit{CrossEntropyLoss} dengan label smoothing $0.1$. Metode \textit{Early Stopping} dengan kesabaran (\textit{patience}) 7 epoch dipasang untuk menghentikan proses latih saat \textit{validation loss} jenuh.

\subsection{Evaluasi Kuantitatif}
Proses pelatihan terhenti otomatis pada epoch ke-19. Hasil akhir evaluasi kinerja model pada data uji disajikan pada tabel berikut:

\begin{table}[H]
\centering
\caption{Metrik Evaluasi Akhir Model Klasifikasi}
\label{tab:metrics}
\begin{tabularx}{\textwidth}{l X c}
\toprule
\rowcolor{tblhead}
\textbf{Nama Metrik} & \textbf{Deskripsi} & \textbf{Nilai Akhir} \\
\midrule
Akurasi Uji (Test Accuracy) & Persentase ketepatan klasifikasi pada data uji & 84.17\% (0.8417) \\
Akurasi Validasi (Val Accuracy) & Akurasi terbaik pada data validasi & 84.13\% (0.8413) \\
Validation Loss Terbaik & Nilai kerugian terendah pada data validasi & 0.8122 \\
Total Kelas Serangga & Jumlah kategori serangga yang dikenali & 118 Kelas \\
Epoch Pelatihan & Jumlah epoch hingga pelatuk early stopping & 19 Epoch \\
\bottomrule
\end{tabularx}
\end{table}

Berdasarkan grafik kurva \textit{loss} dan akurasi (lihat bagian evaluasi visual di notebook), terlihat celah yang sangat sempit antara performa validasi dan latihan, mengonfirmasi keberhasilan strategi augmentasi dan \textit{early stopping} dalam menekan \textit{overfitting}.

%================================================================
% BAB 3: PENGEMBANGAN BACKEND & INTEGRASI GEMINI AI
%================================================================
\newpage
\section{Pengembangan Backend \& Integrasi Gemini AI}

\subsection{Desain Arsitektur API}
Komponen server utama dibangun menggunakan kerangka kerja (framework) \textbf{FastAPI} yang menawarkan kecepatan pemrosesan tinggi berbasis standar kode asinkron (\textit{async/await}). API ini melayani dua rute (routing) utama:
\begin{itemize}
    \item \texttt{GET /health}: Digunakan untuk memantau status keaktifan backend secara berkala, memastikan model klasifikasi PyTorch telah dimuat dengan benar ke memori, dan mengonfirmasi validitas API Key Gemini.
    \item \texttt{POST /analyze}: Rute utama yang menerima berkas gambar dari klien, melakukan prapemrosesan citra instan, memanggil model PyTorch untuk menghitung probabilitas top-3, dan meneruskan prediksi terbaik ke layanan Gemini.
\end{itemize}

\subsection{Integrasi Gemini 2.5 Flash Lite}
Layanan wawasan spesies serangga memanfaatkan SDK \texttt{google-genai} untuk mengakses model \textit{Gemini 2.5 Flash Lite}. Dua fitur utama AI Studio yang diintegrasikan adalah:
\begin{enumerate}
    \item \textbf{ThinkingConfig}: Dikonfigurasi dengan \textit{thinking budget} sebesar 2048 token untuk memaksa model melakukan penalaran internal (chain-of-thought) sebelum memformulasikan teks luaran. Hal ini krusial untuk menjaga ketepatan taksonomi serangga.
    \item \textbf{Google Search Grounding}: Diaktifkan sebagai *tool* eksternal. Fitur ini memungkinkan model Gemini mencari informasi biologi serangga terkini secara langsung di internet, mencegah halusinasi data ilmiah.
\end{enumerate}

\subsection{Mekanisme Ketahanan Fallback}
Sebagai langkah mitigasi jika batas penggunaan API tingkat gratis tercapai (error status 503), backend dilengkapi blok penanganan pengecualian (\textit{exception handling}) yang terhubung dengan \textit{database offline lokal}. Ketika rute utama gagal terhubung ke Google AI Studio, server tidak akan mogok, melainkan mengembalikan data taksonomi statis lebah madu (\textit{apis mellifera}), nyamuk malaria (\textit{anopheles}), dan ordo kupu-kupu ngengat (\textit{lepidoptera}) beserta pesan peringatan keterbatasan layanan.

%================================================================
% BAB 4: IMPLEMENTASI FRONTEND & DESAIN ANTARMUKA
%================================================================
\newpage
\section{Implementasi Frontend \& Desain Antarmuka}

\subsection{Desain Visual}
Antarmuka pengguna dikembangkan dari nol menggunakan \textbf{Next.js} dan \textbf{Tailwind CSS}. Desain visual mengadopsi tema gelap premium (*Dark Violet Theme*) dengan warna latar belakang hitam pekat dan aksen gradasi ungu-ke-merah muda (*violet-to-fuchsia*) yang hidup, menghindari skema warna hijau kuno bawaan template proyek asli.

\subsection{Sistem Komponen & Interaksi Animasi}
Seluruh elemen visual didesain interaktif dengan mikro-animasi dinamis:
\begin{itemize}
    \item \textbf{Entrance Animations}: Komponen navbar, header, kartu kontrol, dan kartu hasil masuk ke layar secara bertahap menggunakan animasi geser (\textit{slide-up/slide-in}) teratur yang dikonfigurasi melalui waktu tunda CSS (\textit{delay-100} hingga \textit{delay-700}).
    \item \textbf{Interactive Dropzone}: Area unggah berkas bertindak sebagai kontainer *drag-and-drop* aktif yang mendeteksi masukan citra, membesar secara visual, dan merubah warna batas garis menjadi violet berpendar saat pengguna menyeret citra di atasnya.
    \item \textbf{Animated Counter & Progress Bar}: Persentase keyakinan model memanjang secara bertahap dan angka persen dihitung maju (\textit{count-up}) secara dinamis demi kenyamanan visual pengguna.
    \item \textbf{Loading State}: Efek animasi titik membal (\textit{bounce-dots}) bergantian dipasang di dalam tombol analisis saat proses asinkron inferensi sedang berjalan.
    \item \textbf{Markdown Rendering}: Konten wawasan biologi yang dikirimkan Gemini dirender rapi menggunakan pustaka \texttt{react-markdown} dan plugin \texttt{@tailwindcss/typography} (\texttt{prose}) yang telah disesuaikan warnanya dengan aksen violet.
\end{itemize}

%================================================================
% BAB 5: TANGKAPAN LAYAR & ANALISIS IMPLEMENTASI
%================================================================
\newpage
\section{Tangkapan Layar \& Analisis Implementasi}

\subsection{Petunjuk Penempatan Gambar (Panduan Screenshot)}
Untuk menyelesaikan dokumen ini, praktikan wajib mengambil tangkapan layar (screenshot) aplikasi saat dijalankan secara lokal di peramban (browser) dan menaruh file gambarnya di dalam folder proyek ini dengan rincian berikut:
\begin{enumerate}
    \item \textbf{screenshot\_home.png}: Ambil gambar tampilan penuh halaman web sebelum ada gambar yang diunggah. Tunjukkan layout bersih *Dark Violet Theme* beserta navbar dan stat bar. Taruh file gambar di \texttt{/home/zeyn/Documents/Final-ML-Lab-2026/screenshot\_home.png}.
    \item \textbf{screenshot\_result.png}: Unggah salah satu gambar uji serangga (misal lebah/kupu-kupu) dan tekan "Analisis Serangga". Ambil tangkapan layar setelah hasil klasifikasi top-3 dan deskripsi Gemini AI termuat sempurna. Taruh file gambar di \texttt{/home/zeyn/Documents/Final-ML-Lab-2026/screenshot\_result.png}.
\end{enumerate}

\subsection{Visualisasi Antarmuka Aplikasi}

\begin{figure}[H]
    \centering
    \includegraphics[width=0.85\textwidth]{Screenshot From 2026-06-04 19-51-06.png}
    \caption{Tangkapan layar halaman utama aplikasi LensArthropoda sebelum pengunggahan gambar.}
    \label{fig:app_home}
\end{figure}

\begin{figure}[H]
    \centering
    \includegraphics[width=0.85\textwidth]{Screenshot From 2026-06-04 19-55-50.png}
    \caption{Hasil klasifikasi spesies serangga (top-3 probabilitas) dan wawasan biologi dari Gemini AI.}
    \label{fig:app_result}
\end{figure}

\subsection{Analisis Respon Gemini AI}
Berdasarkan uji coba langsung terhadap luaran teks Gemini AI, respon yang dihasilkan terbukti memiliki relevansi yang sangat tinggi terhadap gambar masukan. Sebagai contoh, saat model \textit{PyTorch} mendeteksi kelas \textit{apis\_mellifera} dengan probabilitas tertinggi, Gemini memformulasikan struktur klasifikasi biologi ilmiah yang sah dan memaparkan fakta fungsional penting lebah madu sebagai penyerbuk ekosistem utama. Penggunaan pencarian Google (\textit{Google Search Grounding}) menjamin akurasi nama genus dan famili tetap sinkron dengan taksonomi biologi terbaru tanpa mengalami halusinasi teks.

%================================================================
% BAB 6: SARAN PENGEMBANGAN LANJUTAN & KESIMPULAN
%================================================================
\newpage
\section{Saran Pengembangan & Kesimpulan}

\subsection{Saran Pengembangan Model Lanjutan}
Beberapa arah perbaikan masa depan yang diidentifikasi dari keterbatasan sistem saat ini antara lain:
\begin{enumerate}
    \item \textbf{Penanganan Masalah Imbalanced Dataset} \\
    Dataset serangga ini memiliki disparitas jumlah gambar yang besar antar-kelas. Pemanfaatan \textit{Weighted Random Sampler} pada pemuatan data latih (\textit{DataLoader}) atau fungsi kerugian khusus seperti \textit{Focal Loss} sangat direkomendasikan untuk menaikkan performa pengenalan pada spesies serangga langka.
    
    \item \textbf{Migrasi ke Arsitektur Deteksi Objek Lokalisasi} \\
    Sistem saat ini menggunakan klasifikasi tingkat gambar. Di lingkungan pertanian riil, beberapa spesies serangga seringkali muncul bersamaan pada satu tanaman. Migrasi ke deteksi objek berbasis \textit{YOLOv10} akan memampukan sistem mendeteksi koordinat letak serangga (*bounding box*) secara simultan.
    
    \item \textbf{Optimasi Inferensi Melalui Kuantisasi Model} \\
    Model EfficientNet-B3 dalam bentuk berkas \textit{TorchScript} masih memiliki ukuran yang lumayan besar untuk aplikasi web lokal. Penerapan kuantisasi pasca-latih (\textit{Post-Training Quantization} ke format INT8) akan memotong ukuran model hingga 75\% dan melipatgandakan kecepatan inferensi di tingkat backend CPU.
\end{enumerate}

\subsection{Kesimpulan}
Sistem otomatisasi identifikasi serangga pertanian \textbf{LensArthropoda} telah sukses diimplementasikan. Melalui integrasi transfer learning model \textit{EfficientNet-B3} (akurasi uji 84.17\%) dan kepintaran Gemini 2.5 Flash Lite, sistem terbukti andal dalam menyajikan hasil prediksi nama kelas serangga beserta penjabaran biologi yang informatif secara asinkron dengan pengalaman UI yang interaktif dan dinamis bagi pengguna.

\end{document}
